% This is a file that contains the tracking of the activities related to the project during the second week of March 2025.

\documentclass[12pt]{article}

\usepackage[a4paper, margin=1in]{geometry}
\usepackage{graphicx}
\usepackage[backend=biber,style=numeric]{biblatex}
\usepackage[most]{tcolorbox}

\addbibresource{references.bib}

\begin{document}

\author{Joan Ronquillo}

\title{Project Progress - Week 2 of March 2025}
\maketitle

This is a file that contains the tracking of the activities related to the OFC project during the second week of March 2025.
The activities are divided into groups and a summary of the progress is provided.

\section{Initial Status of the Project}
The project has made significant progress up until the second week of March 2025. The key achievements include:
\begin{itemize}
    \item The management of Python modules and compilation of LaTeX projects has been learned.
    \item A first approach to C++ modular development of the functionalities has been initiated.
    \item The reading of Draine's paper on radiative corrections in polarizability \cite{Draine1988ApJ...333..848D} has been initiated.
    \item The difficulties in implementing C++ modules have been discussed with Raúl. And intuition on how to proceed has been provided.
    \item The need to use the Discrete Dipole Approximation (DDA) for the calculation of optical forces for particles of size comparable to the wavelength has been discussed.
\end{itemize}
So far, the project is developing in some areas:

\begin{itemize}
    \item The development of the algorithm to obtain the polarizability of the particles.
    \item The modular development, including makefile documentation and Raul's cpp-template-project repository and notes analysis.
    \item The development of a calculation method for Optical Forces given a polarizability, a field, and a field gradient.
\end{itemize}

\section{Potential Tasks for the Week}

\end{document}
